\documentclass[11pt,a4paper]{article}

\usepackage[LGR,T1]{fontenc}
\usepackage[utf8]{inputenc}
% Transliteraciones con macron+agudo (megalḗtōr, kynṓpidos...)
\DeclareUnicodeCharacter{1E17}{\'{\=e}}
\DeclareUnicodeCharacter{1E53}{\'{\=o}}
\usepackage[greek.ancient,spanish]{babel}
\usepackage[a4paper,margin=2.8cm]{geometry}
\usepackage{microtype}
\usepackage[hidelinks]{hyperref}

% Griego politónico escrito directamente en UTF-8:
% \gr{...} para palabras sueltas; bloques con \foreignlanguage{greek}{...}
\newcommand{\gr}[1]{\foreignlanguage{greek}{#1}}

\hypersetup{
  pdftitle={Leer la Odisea en tiempos iletrados. Tercera parte: Helena}
}

\setlength{\parindent}{1.25em}
\setlength{\parskip}{0.35em}
\raggedbottom

\title{Leer la \emph{Odisea} en tiempos iletrados \\ Tercera parte \\ Helena}
\author{}
\date{}

\begin{document}

\maketitle

Cuando nació mi hija mayor, nos debatíamos entre llamarla Irene o Helena.

Los dos son nombres griegos, los dos son preciosos. \gr{Εἰρήνη} (eirēnē) nombra a una de las tres Horas, hija de Zeus y Temis, junto con Dice (Justicia) y Eunomía (Buen Gobierno). Y su traducción, además, no podía ser más hermosa. Irene significa paz.

El origen del nombre «Helena» estaba menos claro. Es bastante probable que proceda del protohelénico \emph{*Hwelénā}, derivado de \emph{*Swelénā}, nombre de una divinidad solar formado sobre la raíz indoeuropea \emph{*swel-} «brillar, arder, calentar». Helena sería así la «Señora de la Luz».

No era fácil decidir entre la diosa de la paz y la diosa de la luz, pero al final acabamos por elegir Irene, escogiendo paz en lugar de guerra. La guerra de Troya, la guerra de todas las guerras, la guerra más cantada y repetida de la literatura, que no habría ocurrido si Helena no se hubiera fugado con Paris.

¿Pero se fugó? En la \emph{Ilíada} (III, 443-446), Paris le dice:

\begin{verse}
\foreignlanguage{greek}{οὐδ᾽ ὅτε σε πρῶτον Λακεδαίμονος ἐξ ἐρατεινῆς\\
ἔπλεον ἁρπάξας ἐν ποντοπόροισι νέεσσι,\\
νήσῳ δ᾽ ἐν Κραναῇ ἐμίγην φιλότητι καὶ εὐνῇ,\\
ὥς σεο νῦν ἔραμαι…}
\end{verse}

\begin{verse}
ni cuando por primera vez, desde la amable Lacedemonia,\\
navegué habiéndote arrebatado en las naves surcamares,\\
y en la isla Cránae me uní contigo en amor y lecho,\\
te deseé como ahora te deseo…
\end{verse}

El verbo \gr{ἁρπάζω} (harpázō) significa literalmente «arrebatar», así como «arrancar, saquear». De la misma raíz vienen las \gr{Ἅρπυιαι}, las Harpías ---«las arrebatadoras»--- y, quizá, por vía latina, rapiña y rapto.

Así que una posible interpretación, que aparecerá más tarde en la literatura, es que Helena no fue seducida, sino raptada contra su voluntad. El propio Homero mete a los dioses por medio para justificar la infidelidad de la reina. En \emph{Ilíada} III, Afrodita la empuja con amenazas al lecho de Paris y Helena se resiste de palabra e insulta a la diosa, aunque finalmente obedece.

Ahí no acaba la cosa. La famosa leyenda del juicio de Paris ---apenas mencionada por Homero y debida en su mayor parte a mitógrafos muy posteriores--- nos dice que
en las bodas de Tetis y Peleo, progenitores del divino Aquiles, \gr{Ἔρις} (Eris), la Discordia, a la que nadie ha invitado, arroja una manzana de oro con la inscripción \gr{τῇ καλλίστῃ} (tē kallistē), «para la más hermosa». Hera, Atenea y Afrodita, las olímpicas principales, poco humildes ellas, la reclaman. Zeus, que no quiere líos con las mujeres de su familia, se niega a arbitrar y envía a Hermes con las tres al monte Ida, donde Paris apacienta rebaños, sin saber aún que es hijo de Príamo.

Las olímpicas no se andan por las ramas e intentan sobornarle, cada una a su manera. Hera le tienta con el poder, proponiéndole el dominio de Asia. Atenea, por su parte, le promete la victoria en la guerra y la sabiduría. Pero Afrodita le ofrece a la mujer más hermosa del mundo y el pardillo de Paris ni se lo piensa. Afrodita, ¡ay!, olvida explicarle que la hermosa en cuestión, Helena, está casada con el rey Menelao de Esparta.

Los griegos, todo hay que decirlo, son tremendos. Lo primero, el juicio de Paris, de juicio, no tiene nada. Las tres diosas intentan comprar a un adolescente despistado que escoge, como era de esperar ---está en la edad de tener las hormonas altas---, el deseo sobre el poder o la inteligencia, atributos ambos que le tienen sin cuidado a su tierna edad. Lo segundo, una vez que el juicio se resuelve a favor de Afrodita, la cosa no se queda ahí y las otras dos deidades van a conspirar sin vergüenza para arruinar al pobre príncipe ---y de paso a toda su gente--- por una simple cuestión de despecho. Toda la guerra de Troya, según se mire, se debe a una pelea de familia entre las diosas.

El gran Konstantino Kavafis ilustró en muchos de sus poemas la amoralidad de las deidades griegas y, en concreto, nos ofrece una lacerante reflexión sobre las bodas de Tetis y Peleo, a las que acudieron juntos (pero no revueltos) humanos e inmortales:\footnote{El poema se llama «Deslealtad» y cito la magnífica traducción de José María Álvarez en Hiperión, 1976.}


\begin{verse}
En las bodas de Tetis y Peleo,\\
Apolo se levantó en la espléndida mesa\\
del banquete, y llamó bienaventurados a los esposos\\
por el fruto que daría aquella unión.

Dijo: No habrá de ser tocado por la enfermedad\\
y gozará de una larga vida. ---Cuando oyó esto\\
hondamente alegróse Tetis, porque las palabras\\
de Apolo, experto en profecías,\\
le daban seguridad para su hijo.\\
Y cuando Aquiles creciera, siendo\\
gloria de Tesalia su belleza,\\
Tetis recordaba las palabras del dios.

Pero un día llegaron unos ancianos con noticias\\
que narraban la muerte de Aquiles en Troya.\\
Y Tetis desgarró sus purpúreas vestiduras,\\
se quitó y arrojó lejos\\
contra el suelo sus brazaletes y sus anillos.\\
Y en su dolor recordó el pasado;\\
y preguntó qué había hecho el sabio Apolo,\\
dónde estaba el poeta que en el banquete tan elocuentemente profetizase,\\
mientras arrebataban a su hijo en la flor de la edad.

Entonces los ancianos le respondieron que fue el mismo\\
Apolo quien descendiera a Troya, y junto a los troyanos había matado a Aquiles.
\end{verse}


A lo largo de la \emph{Odisea} de Nolan, la única diosa que aparece, Atenea, interpretada por una compungida Zendaya vestida de blancos tules, balbucea unas cuantas frases \emph{new age}, explicándole a Ulises que los dioses, aunque aparentemente invisibles, están por todas partes. Hay que decir que el héroe lleva la misma empanada mental que el director: este ni se atreve a sacar a los dioses de cuerpo presente ni consigue quitarlos del todo de en medio, pues Poseidón y Calipso siguen implícitamente ahí, y el rol de esta última como bella pescadora o hermosa traficante de drogas no es muy convincente si se le restan atributos divinos.

No me extraña que Nolan no quiera sacar a los dioses, para que la película no se le deslice del lado de \emph{Harry Potter} o \emph{El Señor de los Anillos} y prefiera salir del paso con cuatro menciones indirectas y una proyección psicológica del síndrome postraumático del guerrero. ¡Pero qué gran oportunidad perdida! Los dioses no son solo amorales, sino positivamente nefastos. Cada vez que enredan en los asuntos de los hombres, estos últimos acaban mal, mientras que los olímpicos se van siempre de rositas. La \emph{Odisea}, de hecho, puede leerse como el empeño de un hombre en regresar al hogar en contra de la voluntad de un dios, Poseidón. Si este no acaba por aniquilarlo, como a sus compañeros, es porque cuenta con la protección de otros dioses, que si bien lo ayudan, tampoco evitan que las pase canutas. Odiseo, como el resto de los héroes clásicos, no es otra cosa que un juguete en manos de unas deidades corruptas, saturnales, arrogantes, envidiosas y más a menudo malvadas que bondadosas.

De vuelta a Helena, hay una versión muy interesante de lo que sucede después del juicio de Paris.\footnote{Debida a Estesícoro de Hímera (c. 630-555 a.C.).} Hera, que tiene muy mal perder, conjura un \gr{εἴδωλον} (eídōlon), es decir, una copia de Helena, una especie de androide fabricada a partir del aire, que es la que se lleva Paris a Troya. A la verdadera, Hermes la teletransporta a Egipto, con el propósito de que su matrimonio con Menelao quede intacto. Zeus, que es perfectamente consciente de que toda la maniobra va a causar una catástrofe, lo consiente porque está convencido de que en la Tierra sobran hombres y de paso quiere aprovechar para que Aquiles gane gloria. Todo, como puedes ver, ojiclara lectora, bienagudo lector, muy edificante. No es de extrañar que cuando Epicuro comprende que los dioses, si existen, no se ocupan de los hombres, le parezca la mejor de las noticias.

Volvamos a la historia oficial, que, por cierto, tiene miga. Para empezar, en la \emph{Ilíada} no se resuelve qué pasa con Helena (entre otras cosas porque la \emph{Ilíada} termina con los funerales de Héctor y Troya todavía no ha caído). En la \emph{Odisea}, Helena aparece junto a su esposo Menelao cuando Telémaco viaja a Esparta buscando noticias de su padre. En la película de Nolan, la reina aparece con un tremendo tajo en el rostro, que Menelao deja muy claro que es cosa suya. Leamos a Homero en la excelsa traducción de Tobal. Helena aparece en el canto IV, cuando Telémaco y Pisístrato, el hijo de Néstor (el Nestórida), ya se han sentado a la mesa del rey espartano.

\begin{verse}
Su pensamiento y su pecho rumiando en esas estaban,\\
cuando salió de su alcoba techalta y dulcearomada\\
de Ártemis ruecadeoro Helena tal parhermana.\\
Y le arrimó a ella enseguida sillón bienlabrado Adarasta,\\
y Alcipe le echa a sus pies tapete de blanda lana,\\
y una canasta trae Filo de plata que era una dádiva\\
de Alcandra, esposa de Pólibo, rey de la Teba egipciana,\\
donde riquezas sin cuento allí los palacios guardan.
\end{verse}

Tal es la aparición de Helena en la \emph{Odisea}. Homero nos dice que es \gr{ἐϊκυῖα}, «semejante a» Ártemis, semejanza que Tobal captura con el neologismo parhermana. Ártemis, la de la rueca de oro, es la hermana gemela de Apolo, diosa de la caza, de los animales salvajes, de los espacios no cultivados y ---Homero no da puntada sin hilo--- de la castidad. Tiene bemoles que en la entrada en escena de la adúltera más famosa de la literatura, el Poeta la equipare a la diosa más casta del Olimpo. Nada más aparecer, Adarasta y Alcipe (dos esclavas) le arriman un sillón y le echan un tapete de lana en los pies ---palacio real o no, la calefacción central no funcionaba aún en la época---. La tercera sirvienta, Filo, le trae un regalo de Alcandra, reina de Tebas, en Egipto. La conexión con el misterioso país del Nilo no es baladí, como veremos enseguida. En todo caso, la puesta en escena nos habla de una reina mimada y en pleno esplendor.

Una vez que se ha acomodado, Helena se dirige a su marido:

\begin{verse}
¿Sabes tú ya, celestial Menelao, de ser quiénes se honran\\
estos dos hombres que aquí a nuestra casa llegaron ahora?\\
¿Me engañaré, o digo verdad? Soy de un pálpito boca;\\
nunca yo he visto, te digo, persona así igual a persona,\\
no, ni en mujer ni en varón, que es mirarlo y el pasmo me arroba:\\
del de Odiseo grancorazón siento que este es la copia,\\
sí, de Telémaco, apenas nacido y sin padre en la hora en\\
que por mi culpa, perra de mí, los aqueos a Troya\\
fuisteis moviendo una guerra valiente allí y animosa.
\end{verse}

Es Helena quien reconoce a Telémaco, llamándolo hijo de Odiseo grancorazón, otro de los precisos hallazgos de Juanma, que traduce «literalmente» \gr{μεγαλήτορος} (megalḗtoros), genitivo de \gr{μεγαλήτωρ} (megalḗtōr). El prefijo \gr{μέγας} (mégas), de uso corriente en castellano, significa grande o enorme, mientras que \gr{ἦτορ} (ḗtor) es «corazón» en la acepción de sede de ánimo y vida ---el órgano físico es \gr{καρδία} (kardía)---.


Pero además, Helena se increpa a sí misma. En la traducción al castellano, dice: «Por mi culpa, perra de mí». Aquí, excepcionalmente, Tobal se aparta de su propia norma. \gr{κυνώπιδος} (kynṓpidos) es el genitivo de \gr{κυνῶπις} (kynôpis), palabra compuesta por \gr{κύων} (kýōn), «perro», y \gr{ὤψ} (ṓps), «cara, ojo». Un posible neologismo a cuento habría sido «caraperra» o incluso «ojiperra». También es interesante caer en la cuenta de que Aquiles usa un término similar para llamar «caradura» a Agamenón en la bronca entre ambos que abre la \emph{Ilíada} ---\gr{κυνὸς ὄμματ' ἔχων} (kynòs ómmat' échōn), «con ojos de perro»---. En resumen, no queda del todo claro si Helena se está llamando a sí misma «perra» (un epíteto que tienta la lectura de «prostituta») o, simplemente, «descarada». Lo que sí está claro es que lo dice tan campante, como quien confiesa que se ha fumado un cigarrillo a escondidas. La Helena de Homero ni tiene miedo, ni se humilla, ni, desde luego, lleva un tajo en la cara.

Menelao le contesta que él también ha reparado en el parentesco y la evocación de Odiseo, y Pisístrato tercia para confirmar que, en efecto, su compañero es Telémaco. Menelao se arranca entonces a contar el amor que siente por su ausente amigo:

\begin{verse}
Ay, que a mi casa ha llegado el hijo de un hombre que es mucho\\
más que un amigo: ¡tan siempre en la brega por mí y en apuro!\\
Me prometí que a su vuelta a amarlo iba más que a ninguno\\
de los argivos, si en lévedas naves sobre olas y tumbos\\
volver nos daba a los dos Olimpio Zeus miramundo.
\end{verse}

Menelao adora a Odiseo, pero, si las cosas hubieran salido como él quería, la relación asimétrica entre ambos bien podría haber acabado a tortazos. Los bromances tienen su peligro, sobre todo cuando uno de los bros es rey de la poderosa Esparta y el otro el reyezuelo de un árido islote.

\begin{verse}
Que una ciudad yo aquí en Argos le diera y alcázar robusto,\\
para que de Ítaca aquel con su hijo viniese y sus súbditos\\
y con sus bienes, vaciando para ellos de gente algún burgo\\
de los que están a redor y yo señoreo y conduzco.\\
Fundidos una y mil veces, estando él aquí, los dos juntos,\\
ni el buen amor ni el contento tendrían divorcio alguno\\
hasta que nos envolviera la muerte negra en su humo.\\
Y eso con celos un dios debió de mirarlo, seguro, y\\
que solo aquel sin regreso quedara, infeliz, dispuso.
\end{verse}

El caso es que, con la invocación de Menelao, les entran a todos ganas de llorar.
Hay que decir que los aqueos, cuando no están descabezándose los unos a los otros, son muy llorones. Nadie les había explicado todavía que llorar era cosa de mujeres y no se cortaban un pelo.

\begin{verse}
Tal dijo, y les despertó a todos ansia de llanto.\\
De Zeus por padre nacida, Helena ya está llorando,\\
y está llorando Telémaco, y llora también Menelao\\
y ni el de Néstor del lloro sus ojos llega a apartarlos,\\
que es que a su pecho le vino su Antíloco tan sin fallo,\\
a quien de Aurora luciente dio muerte el hijo preclaro;
\end{verse}

Pisístrato se apunta al lloriqueo general acordándose de su hermano Antíloco, muerto en Troya; Menelao se emociona todavía más, pero en esto se acuerda de que si siguen así se va a enfriar la sopa:

\begin{verse}
A ese llorar que nos vino cerrémosle ahora la puerta\\
para pensar en la cena de nuevo, y que el agua nos viertan\\
ya por las manos. Que tiempo tendremos desde que amanezca\\
para contarnos historias Telémaco y yo según vengan.
\end{verse}

Y en ese momento, Homero nos demuestra que Helena es mucho más que la adúltera maltratada que nos pinta Nolan. De hecho, es una hechicera que poco tiene que envidiarle a Circe:


\begin{verse}
He aquí que piensa otra cosa Helena, de Zeus la nacida:\\
echó un remedio de pronto al vino del que bebían\\
contra amarguras y duelos, desacordador de las cuitas.\\
Quien lo tragaba, una vez que en la jarra mezclado se había,\\
lágrima por las mejillas no más echaba ese día,\\
no, ni aunque se le murieran la madre y el padre, o que a pica\\
viera que ante él a un hermano o a un hijo le siegan la vida\\
y con sus ojos lo viera. Remedios como este tenía\\
la hija de Zeus, benignos, de muy buena sabiduría, y\\
fueron merced de la esposa de Ton, Polidamna la egipcia,\\
que allí la tierra feraz mucho gusta de dar medicinas,\\
muchas en mezcla benignas, y en mezcla muchas mortíferas.
\end{verse}

El mejunje es tan potente que el que lo bebe no puede echar más lágrimas (es decir, entristecerse) durante un día aunque le maten a un hijo ante sus propios ojos. La imagen pone los pelos de punta: el filtro que tiene Helena deja al fentanilo a la altura del betún y convierte, de hecho, a un hombre en una piedra durante un rato. Homero nos informa de que «Remedios como este tenía la hija de Zeus, benignos, de muy buena sabiduría», aunque lo de «benignos» lo pongo en duda. Lo que está claro es que Helena dispone de una farmacia considerable, importada del mismo Egipto (donde quizás pasó los diez años que duró la guerra de Troya mientras los dioses les tomaban el pelo a los humanos haciendo que se mataran los unos a los otros por un \emph{eídōlon} que la imitaba). Conoce «medicinas, muchas en mezcla benignas, y en mezcla muchas mortíferas». Así que si no ha envenenado a Menelao, es porque no le ha dado la real (nunca mejor dicho) gana.

Qué quiere exactamente Helena es algo que la \emph{Odisea} nunca deja claro. Un poco más adelante, la reina cuenta cómo Odiseo se infiltra en Troya, disfrazado de mendigo:

\begin{verse}
Hijo de Atreo, feliz Menelao, y vosotros, retoños\\
también de hombres de pro, como el dios ora al uno ora al otro,\\
Zeus, lo bueno y lo malo lo da, porque puédelo todo,\\
banquetead bien ahora, sentados en casa, y buen gozo\\
de las palabras tened: lo que os voy a contar viene a tono.\\
Todas no puedo narrar, ni siquiera a nombrarlas me pongo,\\
cuantas proezas me sé de Odiseo biencordianimoso;\\
ved solo qué hizo, valiente, metiéndose en boca del lobo\\
por esos pagos de Troya, de aqueos peste y ahogo.\\
Tras lacerarse a sí mismo con golpes muy vergonzosos\\
y echarse, igual que un criado, andrajos por sobre los hombros,\\
en la ciudad calliclara de aquellos se entró, de los otros;\\
y otro él también parecía envuelto en ese rebozo,\\
un de pedir: nunca en nave de aqueos se vio en ese cómo.
\end{verse}

La escena, que Nolan también omite, muestra a la vez la valentía de Odiseo ---hay que tener agallas para colarse de incógnito en la ciudad enemiga--- y su astucia, ya que solo a él, entre toda aquella banda de soldados fornidos y quizás algo espesos de mollera, se le ocurre disfrazarse de mendigo. Para hacer su disfraz realista, además, Odiseo se golpea a sí mismo «con golpes muy vergonzosos» y se echa sobre los hombros andrajos de criado. El ardid, por supuesto, adelanta la treta a la que el héroe va a recurrir cuando regrese a Ítaca y, otra vez, finja ser un mendigo. Y si me apuras, atenta lectora, perspicaz lector, le da contexto a la caracterización de sí mismo que más tarde Odiseo hará con el cíclope, poniéndose el nombre de \gr{Οὖτις} (outis), es decir, «Nadie». El truco más famoso de la literatura ---que, increíblemente, Nolan también pasa por alto--- le sirve para confundir a los otros cíclopes, como veremos más adelante, y se convierte en su alter ego. Y si Ulises es Nadie, parece apropiado que se vista como tal: un mendigo, un andrajoso, el más humilde de los sirvientes.

Pero Helena lo reconoce. Es más, asegura que el astuto Odiseo, que no se fía ni de su sombra, acaba confiando en ella:

\begin{verse}
Y tal que así, en el alcázar de Troya se entró, y voz ni ojos\\
tuvo de nadie; que aquel era él supe verlo yo solo,\\
y lo llamé; pero en mañas y trampa volvíame el rostro.\\
Mas cuando lo hube bañado y ungido de oliva, y le pongo\\
ropas encima y le juro con juro bien riguroso\\
que no daré de Odiseo a los troyanos un soplo\\
antes de que haya tocado su tienda y naves a poco,\\
de los aqueos entonces me cuenta cuál el propósito.
\end{verse}

No solo ayuda a Odiseo. Helena afirma alegrarse cuando por fin cae la ciudad, asegura que echa de menos su patria, a su hija y a su esposo, culpa a Afrodita y confiesa un poco de remordimiento.

\begin{verse}
Y de troyanos mató con su bronce afilado buen monto\\
hasta que al fin con los suyos volvió, y con saberes valiosos.\\
Allí a las otras troyanas yo oía llorar; y en el fondo\\
yo estaba alegre, que mi corazón ya ansiaba el retorno,\\
que me cegó a mí Afrodita ---sentía ahora ya reconcomio---\\
cuando de tierra y de patria hasta allí me llevó, y en divorcio\\
de esa mi hija y de alcoba y de un esposo que a otro\\
en nada cede en estampa ni en pensamiento juicioso.
\end{verse}

En resumen, una adúltera arrepentida que quiere volver a casa y a la que, por cierto, no le cabe duda de que así será. Los versos no muestran miedo a la venganza del marido, ni sombra de duda sobre cuál puede ser su destino. Demuestran, eso sí, una sangre fría notable. Helena lleva diez años en la corte del rey Príamo, ha compartido el lecho con dos príncipes, Paris y Deífobo\footnote{Aunque de su unión con Deífobo la \emph{Odisea} solo deje una huella de refilón, su historia, como la de muchos otros actores secundarios de la \emph{Ilíada} y la \emph{Odisea}, se cuenta en los poemas del Ciclo épico que «complementaban» la \emph{Ilíada}. No hay que olvidar que los hechos que Homero relata en esta corresponden a unas pocas semanas del décimo año de la guerra y se omite la mayoría de las historias que, con el paso de los siglos, irían añadiéndose. Así, los \emph{Cantos ciprios} van desde las causas de la guerra hasta el arranque de la \emph{Ilíada}; la \emph{Etiópida} recoge lo que sigue a la muerte de Héctor y en ella aparecen personajes tan inolvidables como Pentesilea, además de relatarse la muerte de Aquiles; la \emph{Ilíada menor} llega hasta el famoso caballo; la \emph{Iliupersis} narra el saqueo de la ciudad y los \emph{Nostoi}, los regresos de los diversos príncipes aqueos. Los hexámetros de esas sagas no se conservan, pero sí las tramas, gracias a un resumen en prosa atribuido a cierto Proclo, transmitido por manuscritos medievales de la propia \emph{Ilíada} y por la \emph{Biblioteca} de Focio (c. 810-893), patriarca de Constantinopla.}, es la causante directa de la destrucción de la ciudad y todo lo que tiene que decir es: «Allí a las otras troyanas yo oía llorar; y en el fondo / yo estaba alegre, que mi corazón ya ansiaba el retorno».

Pero hay más. Homero no solo nos dibuja a una Helena radiante, poderosa y fría como un témpano (es ella la que, con su filtro, corta los lamentos de los invitados y acaba de relatarnos que los llantos de las troyanas le tienen sin cuidado). También nos advierte de que es una narradora de la que no podemos fiarnos. Porque Menelao, su marido, le contesta así:

\begin{verse}
Sí, cuánta ley y razón, mujer, en lo que has bien traído.\\
De hombres heroicos ---de muchos--- pensares y miras he visto,\\
pues que por muchas también de tierras hice camino,\\
pero jamás corazón vi yo con los ojos míos\\
que al de Odiseo biencordianimoso le fuera vecino.\\
Ved solo qué hizo, valiente, arriesgándose dentro metido\\
de aquel pulido caballo, y con él los mejores argivos\\
que a los troyanos parca y matanza a llevarles ya íbamos.
\end{verse}

Menelao insiste en lo mismo que Helena. Ulises no solo es astuto ---inventándose la treta del caballo--- sino valeroso, siendo el primero que se mete dentro, acompañado por el propio Menelao, que tampoco es ningún pusilánime. Ciertamente, todos los asaltantes ---«los mejores argivos»--- se juegan el tipo en el ataque, algo que Nolan capta, en mi opinión, de forma muy solvente\footnote{Entre las críticas que se le han hecho está la de que el caballo no lleve ruedas, pero lo cierto es que Homero no se las pone en ningún momento: todo lo que dice es que el caballo es arrastrado al interior de la ciudad. La tradición posterior, por supuesto, empezando por Virgilio en la \emph{Eneida}, sí se las pone. Anterior a Virgilio tenemos el \emph{pithos} de Míconos, un vaso de relieve circa 670 a.C., que se considera la representación más antigua conservada del caballo, y lleva ruedas visibles, además de escotillas con guerreros asomando. Nada le impedía a Nolan ponerle ruedas al equino, pero la decisión de no ponérselas no tiene nada de malo.}.

Pero ahora viene lo bueno. Refiriéndose a Helena, sigue Menelao:

\begin{verse}
Tú te acercaste allí entonces (te habría un diosdalgo movido\\
por gloria en mano tenderles a los troyanos ahí mismo)\\
e iba guardando tus pasos Deífobo ludivino.\\
Tres vueltas diste tentando a placer el rehueco armadijo,\\
y de los dánaos decías los nombres más distinguidos\\
poniendo voz de su esposa en cada nombre de argivo.
\end{verse}

Menelao afirma literalmente que Helena quiere ayudar a los troyanos, utilizando, una vez más, sus poderes casi mágicos, que le permiten llamar a los aqueos fingiendo la voz de cada una de sus esposas. Y sigue:


\begin{verse}
Sí, y el divino Odiseo y el de Tideo y yo mismo,\\
justo en la panza sentados, llamarnos así te oíamos.\\
Ya no aguantábamos más nosotros dos en delirio\\
por dar un salto y salir o hablarte a ti en un respingo,\\
cuando Odiseo del ansia nos tuvo y barrera nos hizo.
\end{verse}

Nolan escenifica el momento en que los aqueos son casi descubiertos mediante la violencia: los troyanos pinchan el caballo e hieren a uno de los guerreros mientras el populacho exige que se le pegue fuego al invento, lo que a su vez crea un instante de pánico entre los soldados atrapados en el ingenio que consiguen reprimir \emph{in extremis}. El resultado no funciona mal, pero renuncia a una escena mucho más compleja e interesante, en la que Menelao y Diomedes («el de Tideo») están a punto de delatarse y es, una vez más, Odiseo quien salva la situación:

\begin{verse}
Allí en los hijos de aqueos se hacía, en todos, mutismo,\\
que, aunque palabra y respuesta quería dejar ir Anticlo,\\
con ambas manos estaba Odiseo apretándole el pico,\\
manos de fuerza, y sin pausa lo tuvo, de aqueos abrigo,\\
hasta que Palas Atena de allí te llevó al fin consigo.
\end{verse}

Homero no nos ofrece aquí sangre y fuego, sino la confrontación de dos enemigos a cuál más inteligente, Helena y Odiseo ---añadiendo la vuelta de tuerca de que la una ha declarado su amistad por el otro unos versos antes---. Cada uno despliega su propia astucia y el duelo está muy igualado, casi podríamos decir que los aqueos se salvan por la proverbial campana.

¿Y qué ocurre después de que Menelao le llame embustera ---y traidora--- a su esposa delante de sus invitados? Nada de nada. Telémaco, bien porque está agotado, bien porque quiere echar balones fuera, cambia de conversación y pide un catre:

\begin{verse}
Y ya Telémaco ahora, mirándolo, así está diciendo:\\
«Gran Menelao, Atrida celeste, caudillo de ejércitos,\\
triste es, bien triste, que ni eso librarlo del trance más negro\\
pudo, por más que de hierro de sí un corazón iba dentro.\\
Dadnos ya un lecho y que a él se nos lleve, y que en algo contento\\
pueda, adormidos, venirnos al aura del dulce sueño».
\end{verse}

Y es Helena, naturalmente, la que se ocupa de acomodarlos (de paso Homero nos vuelve a recordar que en Esparta refresca bastante):

\begin{verse}
Dijo, y Helena la argiva mandando está ya a las criadas\\
catres armar bajo el soportal, y por bajo unas mantas\\
bellas de grana extender, y por cima unas colchas de lana\\
y unos manteos enteros porque resguardo les hagan.\\
Y con la antorcha en su mano saliendo ahora están de la sala,\\
y les alisan los catres. Y allí ya un heraldo los guiaba.\\
Que allí adurmiéronse ellos, en el portal de la estancia,\\
de Néstor lígrimo el hijo y Telémaco divinarraza.\\
El Atreída en lo hondo de la alta mansión se acostaba,\\
y Helena mantodevuelo con él se encamó, damisanta.
\end{verse}

¿Cómo acaba la noche? Pues como la de un matrimonio cualquiera que, si acaso, ha discutido un poco durante la cena. Menelao se va a dormir y Helena con él. Como si tal cosa. Tobal traduce \gr{δῖα γυναικῶν} con el estupendo neologismo «damisanta», pero una alternativa en prosa sería algo así como «divina entre las mujeres».

De hecho, es evidente que Homero considera divina a Helena. Tanto es así que Menelao se agencia, gracias a ella, un pasaje en primera al Elíseo, es decir, lo más parecido al Paraíso cristiano que los antiguos griegos conocían. Se lo cuenta Proteo, el viejo del mar, en el mismo libro IV que estamos visitando.


\begin{verse}
En cuanto a ti, celical Menelao, voz de dios no refiere\\
que en Argos mueras yegüera y tu sino allí te lo encuentres,\\
no, que a la campa de Elíseo, la tierra final, los sin muerte\\
te llevarán, donde está Radamantis el güero, y es siempre\\
para los hombres muy fácil la vida, que allí nunca hay nieve\\
ni hay un invierno que mucho resulte y tampoco allí llueve,\\
que solo el soplo sonoro del vendaval del poniente\\
manda a los hombres Océano porque con él se refresquen;\\
y es porque tienes a Helena, y yerno de Zeus te sienten.
\end{verse}

En resumidas cuentas, Homero nos presenta a Helena como una deidad más, con impunidad para hacer lo que le dé la gana sin que nadie le pase factura. Y esa impunidad, que es escalofriante, es el elemento central del personaje. Helena es la diosa de la que se sirven otros dioses para sembrar el dolor y la muerte entre los hombres ---por sus propios y banales motivos--- sin que se le exija en ningún momento responsabilidad por ello. Aún más, Menelao hace un magnífico negocio perdonándole sus pecadillos, ya que con ello se garantiza la eternidad en las praderas luminosas del Elíseo, nada que ver con el horrible Hades que el propio Ulises visita en la \emph{Odisea} y del que hablaremos más adelante\footnote{El Hades poco tiene que envidiarle al infierno de Dante, pero no está poblado por las almas de los réprobos, sino que allí va a parar todo el mundo: buenos, malos, valientes, cobardes, sabios y obtusos. El Elíseo, en cambio, no es el destino común, sino un privilegio excepcional que los inmortales conceden a Menelao por ser yerno de Zeus. Ni siquiera Aquiles, hijo de una diosa, tiene carnet de socio; tampoco Agamenón ni el propio Odiseo.}.

Nolan pierde una gran oportunidad por no haber leído con atención el canto IV, o no haberse atrevido a ser fiel a Homero. En lugar de una semidiosa, nos presenta a una mujer maltratada, lo cual no es solo una falsa lectura, sino un cliché.
Hay una tendencia tan aburrida como nefasta en Hollywood y en buena parte de la narrativa actual a zanjarlo todo con un discurso moral que, en último término, premia a los buenos o, como mínimo, castiga a los malos. En Homero no hay moralina que valga. Hay personajes perfectamente buenos que perecen y otros rematadamente malos que se salen con la suya. Los héroes acaban en el Hades y los cuñados de los olímpicos en el Elíseo. Y la adúltera causante de la mayor catástrofe narrada por el mundo antiguo vuelve radiante a su palacio sin que nadie le toque un pelo de la divinal cabeza.



\end{document}
